Hej 
Jeg ønsker hermed at indgive en formel klage over bedømmelsen af min SRP. 
Jeg har modtaget karakteren 02, og efter at have gennemgået både min opgave og mit forsvar er jeg i 
tvivl om, hvorvidt bedømmelsen er retvisende for det faglige niveau i min præstation. 
Under forsvaret oplevede jeg nogle forhold, som efterfølgende har gjort mig usikker på 
bedømmelsesgrundlaget. Blandt andet gik censorens computer tør for strøm under eksaminationen, 
hvilket påvirkede forløbet. Jeg oplevede også flere situationer, hvor censor virkede usikker på 
formuleringen af sine egne spørgsmål, hvilket gjorde det vanskeligt for mig at forstå præcist, hvad der 
blev spurgt ind til. 
Efter voteringen fik jeg oplyst, at min redegørelse og analyse ikke blev vurderet som tilstrækkeligt 
dybdegående. Jeg har dog ikke modtaget en nærmere faglig begrundelse for, hvilke konkrete mangler 
der førte til den samlede karakter på 02. Derfor har jeg haft svært ved at forstå, hvordan bedømmelsen 
er nået frem til netop denne karakter. 
Jeg har efterfølgende kontaktet min lærer og bedt om en skriftlig begrundelse for karakteren, da jeg 
gerne ville forstå vurderingen bedre. Jeg fik imidlertid oplyst, at han først ønskede at tage en mundtlig 
samtale om bedømmelsen efter klagefristens udløb. Dette har gjort det vanskeligt for mig at få en 
nærmere forklaring på bedømmelsen inden for den periode, hvor jeg skal tage stilling til en eventuel 
klage. 
Jeg mener, at min opgave besvarer alle dele af opgaveformuleringen og indeholder både redegørelse, 
arbejdskravsanalyse, analyse, diskussion og konklusion. Derfor oplever jeg, at der er en betydelig forskel 
mellem min egen vurdering af opgavens faglige niveau og den endelige bedømmelse. 
På den baggrund ønsker jeg at klage over bedømmelsen af min SRP og få sagen behandlet. 
Med venlig hilsen 
Sofie Hageskov Kjellegaard 
3.B
